\chapter*{Abstract}
\addcontentsline{toc}{chapter}{Abstract}

Groundwater monitoring networks cover only a fraction of the locations where predictions are needed, and most forecasting systems work only where wells already exist.
Forecasting groundwater levels (GWL) at existing monitoring wells improves considerably with
deep learning. However, predicting GWL at entirely unmonitored locations without any
historical observations is more difficult and receives less attention. Most
existing approaches do not test against large, fixed test sets of spatially separated
wells and therefore cannot "forecast anywhere".

This thesis introduces GIST (\textbf{G}roundwater \textbf{I}nterpolation in \textbf{S}pace
and \textbf{T}ime), a spatiotemporal pipeline combining a temporal and a spatial modeling component,
covering 1{,}040 monitoring wells across Brandenburg, Germany.
A Gated Recurrent Unit (GRU) produces temporal forecasts at monitored wells, which a Gaussian Process
(GP) then interpolates to locations excluded from training to simulate unmonitored locations.

The GRU matches the performance of a much larger Temporal Fusion Transformer baseline while
training roughly 90 times faster, which makes the full experimental setup feasible. One of the main findings is that temporal forecast quality is not the bottleneck
in the spatiotemporal pipeline: interpolating the actual observations at training wells instead
of the temporal predictions does not improve the error of the full spatiotemporal pipeline. The limiting factor is therefore spatial generalization, not temporal forecast quality. Distance to the nearest training 
well is the strongest predictor of per-well error. A jointly trained model combining GRU and GP with a combined loss does not improve overall performance but specifically reduces error for the most spatially isolated wells,
which are the hardest cases for a ``forecast anywhere'' system.
Together, these results demonstrate the feasibility of spatiotemporal GWL forecasting at unmonitored locations and show that further progress mainly depends on improving spatial generalization.
